% !PRINTED 498
\weberpage{498}

If we set
\weq{3}{a_0 = m\,e_0,\quad a_1 = m\,e_1,\quad a_2 = m\,e_2 \ldots,}
then the $e_0, e_1, e_2, \ldots$ are integers without a common divisor;
\weq{4}{E(x,y,z,\ldots) = \sum e\,x^r y^s z^t \ldots}
is a primitive function, and we have
\[ \Phi = m\,E. \]

We have considered these functions $\Phi$ in \S.\,2 of the first volume, and have there proved of them the theorem:

3. \gesperrt{The divisor of a product of two or more integral functions $\Phi$ is equal to the product of the divisors of the individual factors.}

The integral functions in $\Omega$ depend, besides on the variables, on an algebraic number $\Theta$, but otherwise contain only rational numerical coefficients. If we denote such a function by $\varphi(\Theta, x, y, z, \ldots)$, we obtain the \gesperrt{conjugate functions} $\varphi, \varphi_1, \varphi_2, \ldots$ or
\weq{5}{\varphi(\Theta,x,y,z,\ldots),\quad \varphi(\Theta_1,x,y,z,\ldots),\quad \varphi(\Theta_2,x,y,z,\ldots),\ \ldots}
when we substitute for $\Theta$ all the roots of the irreducible equation $f(x)=0$ [\S.\,132, (2)]. These conjugate functions may also in part be equal to one another.

They are all equal to one another when $\varphi$ is a function in $R$; and conversely $\varphi$ is a function in $R$ when the conjugate functions are all equal to one another; for in that case, if we understand by $S(\varphi)$ the sum of the conjugate functions (the trace),
\[ n\varphi = S(\varphi), \]
and $S(\varphi)$ is an integral function whose coefficients are symmetric functions of the $n$ roots $\Theta$, that is, rational numbers.

Among the integral functions in $R$ there belongs also the \gesperrt{norm of} $\varphi$, that is, the product
\weq{6}{N(\varphi) = \varphi\,\varphi_1\,\varphi_2 \ldots,}
for all the coefficients of this function are symmetric functions of the $\Theta$. This function is divisible by $\varphi$, and if we set
\[ N(\varphi) = \varphi\,\varphi' \]
then both $\varphi$ and $\varphi'$ are integral functions in $\Omega$. For $\varphi'$, as the product $\varphi_1\varphi_2\ldots$, is an integral function, and the coeffi-
%CONT>>
